\subsection{从分配块到复用方向}

最初的 Leech 提升保留大型母壳，采用已有496点相容块。
可改变的关键变量是块与尾标签的分配。
25维中，高度为二分之一的纬度允许加入两极点，
无需扩大块就得到197058点。

27维中，十二个标签可以分成四个零和三角形。
将块的复用次数从 $(3,3,2,2,2)$ 改为 $(3,3,3,3)$，
就能多保留一个496点赤道块，得到200540。
容量论证解释了这个点数：每个方向至多复用三次，
等号要求标签装满且每个已用方向都复用三次。

后来的201567点结果使用另一份公开输入。
Lindow 的构造联合优化第一、第二层，得到201566点。
新的见证保留其第一层，把第二层的所属方向从310个增加到311个；
第一层的2090个需删除赤道点的头和168个无须删除的头不变。

扣除自身删除代价后，新增头可能贡献二点或三点，
而第二层每个所属方向贡献一点。
只优化一层的总点数，可能反而选出较差的完整构型。
